\customchapter{ABSTRACT}

\begin{center}
\large \vspace{-1.5cm} \textbf{HyperNCD-Sonata hybrid model for novel class
discovery in the semantic segmentation of architectural heritage point clouds}
\end{center}

This research focuses on recognizing unlabeled architectural elements in point
clouds of built heritage, using the Cathedral of Lima as a case study. Laser scanning records
where matter is, but not which element each point stands for, and that
interpretation is still done by hand. Current segmentation methods
recognize only the categories they were trained on.

We propose a hybrid model that starts from HyperNCD, a novel class discovery
method, and replaces its geometric descriptor, computed with a fixed formula, with
the self-supervised representation of Sonata, used without retraining. The cloud
is prepared with outlier removal (SOR), voxelization and normals. Each point gets a semantic vector, learned by HyperNCD, and a
self-supervised one, and both are combined to cluster the points.

The results show that HyperNCD trained on urban scenes does not transfer to the
cathedral: almost all points (98.4\%) fall into a single output of the model.
Training it on ArCH, an architectural-heritage dataset~\cite{matrone2020benchmark},
reduces that concentration to 44.6\% and raises the mIoU to 35.5\%. When the two
descriptors are combined, the mean mIoU is 52.5\% against 51.4\% without Sonata, a
difference that is not conclusive, and adapting Sonata to the cathedral does not
improve the result. The evidence comes from one scene and one training run, so more
scenes and repetitions are needed.

\vspace{0.5em}

\noindent\begin{minipage}{\linewidth}
\noindent\textbf{Keywords:}\par
\noindent Point clouds; Three-dimensional semantic segmentation; Novel class
discovery; Self-supervised learning; Hypergraphs; Architectural heritage
\end{minipage}
